\section{Kesimpulan}
\label{sec:kesimpulan}

Berdasarkan hasil penelitian yang telah dilakukan, diperoleh beberapa kesimpulan sebagai berikut.

\begin{enumerate}

\item Pada ruang-waktu Schwarzschild, bagian real dan imajiner frekuensi mengalami
penurunan seiring meningkatnya massa lubang hitam, sedangkan pada Schwarzschild Anti-de
Sitter kedua bagian frekuensi justru meningkat seiring meningkatnya jari-jari horizon
peristiwa. Peningkatan momentum sudut menghasilkan kecenderungan yang sama pada kedua
ruang-waktu, yaitu bagian real meningkat sedangkan bagian imajiner menurun, sehingga laju
peluruhan perturbasi melambat. Sebaliknya, peningkatan indeks \textit{overtone}
menghasilkan kecenderungan yang berbeda antara kedua ruang-waktu: pada Schwarzschild,
bagian real menurun sedangkan bagian imajiner meningkat, sedangkan pada Schwarzschild
Anti-de Sitter kedua bagian frekuensi meningkat. Pada penambahan massa medan skalar,
ruang-waktu Schwarzschild meningkatkan frekuensi osilasi namun menurunkan laju
peluruhannya, sedangkan pada Schwarzschild Anti-de Sitter penambahan massa medan
meningkatkan baik frekuensi osilasi maupun laju peluruhannya.

\item Analisis terhadap bagian imajiner frekuensi menunjukkan bahwa seluruh mode yang
berhasil diperoleh memiliki nilai $\omega_I>0$, baik pada ruang-waktu Schwarzschild maupun
Schwarzschild Anti-de Sitter. Hal ini menunjukkan bahwa kedua geometri berada dalam kondisi
stabil terhadap perturbasi medan skalar, baik untuk medan skalar tak bermassa maupun
bermassa. Meskipun sama-sama stabil, karakteristik peluruhan perturbasi pada kedua
ruang-waktu berbeda. Pada ruang-waktu Schwarzschild, peningkatan massa lubang hitam
menyebabkan laju peluruhan perturbasi semakin lambat, sehingga waktu yang dibutuhkan
perturbasi untuk meluruh menuju kesetimbangan menjadi lebih lama, sedangkan pada
ruang-waktu Schwarzschild Anti-de Sitter peningkatan ukuran horizon peristiwa justru
mempercepat peluruhan perturbasi. Dengan demikian, geometri ruang-waktu memberikan
pengaruh yang signifikan terhadap karakteristik spektrum frekuensi kuasinormal sekaligus
terhadap laju relaksasi perturbasi menuju kesetimbangan.

\end{enumerate}